\section{Inverse minor rigidity in structured classes}
\label{sec:inverse}

Let $W$ be an invertible nonnegative column-stochastic $n\times n$
matrix, and write
\[
 B=W^{-1},\qquad
 M_k(B)=\max_{|I|=|J|=k}|\det B_{I,J}|,\qquad
 a(B)=\max\{0,-\min_{i,j}B_{ij}\}.
\]
The general problem asks for a dimension-free bound
$a(B)\leq C(c-1)$ when $M_k(B)\leq c^k$ for every $1\leq k\leq n$,
with $c\geq1$.  We prove such bounds for support-two matrices,
low-order blocks, and symmetric positive-definite matrices.  The
first class yields an all-scale geometric estimate for every
prescribed maximum simplex of an edge-supported body.  The
unrestricted inverse-minor problem is not resolved by these results.

\subsection{Weighted cycles}

\begin{lemma}[Weighted cyclic blocks]
\label{lem:inverse-cycles}
Suppose a column-stochastic block of order $q\geq2$ has the form
\[
 W=(\Id_q+P R_d)D,\qquad
 D=\diag(d_i),\quad R_d=\diag(r_i),\quad
 d_i=\frac1{1+r_i},\quad r_i>0,
\]
where $P$ is a cyclic permutation matrix.  If
$M_1(B)\leq c$ and $M_2(B)\leq c^2$, then
\begin{align}
 a(B)&\leq c^2-1 &&(1\leq c<\sqrt2),\label{eq:inverse-cycle-local}\\
 a(B)&\leq(2+\sqrt2)(c-1)&&(c\geq1).
 \label{eq:inverse-cycle-global}
\end{align}
\end{lemma}

\begin{proof}
Let $R=\prod_i r_i$.  Interchanging the two cyclic matchings by a
row permutation replaces this product by $R^{-1}$.  We may choose
the orientation with $R\leq1$.  In even order, $R=1$ is excluded by
invertibility.  Thus
\[
 L=1-(-1)^qR>0.
\]
The finite geometric identity gives
\begin{equation}
 B=\frac1L D^{-1}
 \bigl[\Id_q-PR_d+(PR_d)^2-\cdots+
                   (-1)^{q-1}(PR_d)^{q-1}\bigr].
 \label{eq:inverse-cycle-expansion}
\end{equation}
In particular,
\begin{equation}
 B_{ii}=\frac{1+r_i}{L},\qquad
 \det B_{\{i,j\},\{i,j\}}=
 \frac{(1+r_i)(1+r_j)}{L}\quad(i\ne j).
 \label{eq:inverse-cycle-two-minor}
\end{equation}
For the second formula, Jacobi's complementary-minor identity
applies.  Directly, when $q>2$, deleting vertices $i,j$ from the
directed cycle leaves an acyclic principal submatrix, triangular
after topological ordering, with determinant the product of its
$d_\ell$.  For $q=2$, the formula is simply
$\det B=1/(d_1d_2L)$.

First suppose every $r_i\leq1$.  If $q$ is odd, then $R\leq r_i$,
so $B_{ii}\geq1$; if $q$ is even, $L\leq1$ gives the same
conclusion.  In row $i$, each absolute inverse entry in
\eqref{eq:inverse-cycle-expansion} is the common factor
$(1+r_i)/L$ times a product of consecutive ratios.  Negative entries
correspond to odd-length paths.  Every such product is at most the
last one-edge factor $r_j$, where $j$ is the predecessor of $i$.
Thus the largest negative magnitude in that row occurs on the
incoming edge and equals
\[
 -B_{ij}=\frac{(1+r_i)r_j}{L}.
\]
The principal two-minor on $i,j$ is
\[
 \frac{(1+r_i)(1+r_j)}L
 =B_{ii}+(-B_{ij})\geq1+(-B_{ij}).
\]
It follows that $a(B)\leq M_2(B)-1\leq c^2-1$.

Next suppose some ratio is at least one, and let $x\geq y$ be the
two largest ratios, with $x\geq1$.  Then $R\leq xy$: if $y\leq1$
all the remaining ratios are at most one, while if $y\geq1$ then
$xy\geq1\geq R$.  In odd order,
\begin{align*}
 (1+x)(1+y)-2(1+R)
 &\geq(1+x)(1+R/x)-2(1+R)\\
 &=(x-1)(1-R/x)\geq0.
\end{align*}
The corresponding two-minor in
\eqref{eq:inverse-cycle-two-minor} is therefore at least two.  In
even order $L\leq1$ and $(1+x)(1+y)\geq2$ give the same result.
Hence this case requires $c\geq\sqrt2$.

Consequently $c<\sqrt2$ forces the first case and proves
\eqref{eq:inverse-cycle-local}.  For $c\geq\sqrt2$, the entry bound
$a(B)\leq c$ yields
\[
 \frac{a(B)}{c-1}\leq\frac c{c-1}\leq2+\sqrt2.
\]
For $1<c<\sqrt2$, the first case gives
$a(B)/(c-1)\leq c+1<1+\sqrt2$.  At $c=1$ its conclusion is
$a(B)=0$.  This proves \eqref{eq:inverse-cycle-global}.
\end{proof}

Independent row and column permutations of $W$ merely permute the
rows and columns of $B$ independently and change signs of its minors.
All these inequalities are therefore invariant under such
permutations.  Zero edge weights may be included by continuity, but
the support decomposition below handles them without treating a
zero-weight edge as part of a genuine cycle.

\subsection{Columns supported on at most two rows}

\begin{theorem}[Support-two inverse-minor rigidity]
\label{thm:inverse-support-two}
Suppose every column of $W$ has at most two nonzero entries.  If
$M_1(B)\leq c$ and $M_2(B)\leq c^2$, then
\begin{align}
 a(B)&\leq c^2-1 &&(1\leq c<\sqrt2),
 \label{eq:inverse-support-local}\\
 a(B)&\leq(2+\sqrt2)(c-1)&&(c\geq1).
 \label{eq:inverse-support-global}
\end{align}
For $n=1$ only the order-one bound is needed.  The constants do not
depend on the matrix order, cycle lengths, or sizes of attached trees.
\end{theorem}

\begin{proof}
Associate a multigraph to $W$.  Its vertices are the row indices;
a column with two positive entries is an edge between its two
supporting rows, and a singleton column is a loop at its supporting
row.  Parallel edges are allowed.  Connected components partition
the rows and columns into matrix blocks.  Each block must be square:
if a component has unequal row and column counts, summing the
component rank bounds, each at most the smaller count, makes the
whole square matrix singular.  Every connected component therefore
has equally many edges and vertices and is unicyclic, including
loops and doubled edges as cycles of lengths one and two.

Remove the trees attached to the unique cycle.  Pair each non-core
vertex $v$ with the edge joining it to its parent $p(v)$ towards
the core.  Order core rows and columns first and tree vertices in
parent-before-child order.  Then
\[
 W=\begin{pmatrix}W_C&F\\0&T\end{pmatrix},
\]
where $T$ is triangular.  The column paired with $v$ has weight
$d_v>0$ at $v$ and $1-d_v$ at its parent.  Its full inverse diagonal
entry is $1/d_v$, so
\begin{equation}
 r_v:=\frac{1-d_v}{d_v}=\frac1{d_v}-1\leq c-1.
 \label{eq:inverse-tree-ratio}
\end{equation}
The core inverse $B_C=W_C^{-1}$ is a submatrix of $B$ and inherits
the order-one and order-two bounds.  A loop core has $B_C=[1]$;
the other cores satisfy Lemma~\ref{lem:inverse-cycles}.

For completeness, let $b_v=B e_v$ denote an inverse column.  The
identity $BW e_v=e_v$ gives the exact recursion
\begin{equation}
 b_v=d_v^{-1}e_v-r_v b_{p(v)}.
 \label{eq:inverse-tree-recursion}
\end{equation}
For a core index $p$, $b_p$ has core part $B_Ce_p$ and zero tree
part.  Repeatedly applying \eqref{eq:inverse-tree-recursion} shows
that every off-diagonal entry in a tree column is a product of at
least one tree ratio, times a core inverse entry or an ancestor
diagonal entry $1/d_u$.  Every such starting factor has absolute
value at most $c$.  Equivalently, the tree block has the form
$(\Id+N)D$ with $N_{p(v),v}=r_v$, and the finite nilpotent expansion
of $(\Id+N)^{-1}$ gives these unique path products.

If $c<\sqrt2$, then $0\leq r_v\leq c-1<1$.  Every tree-column
off-diagonal magnitude is therefore at most
$c(c-1)\leq c^2-1$.  Tree diagonal entries are positive, and core
negative entries satisfy the same $c^2-1$ bound by the cyclic
lemma.  This proves \eqref{eq:inverse-support-local}.  For
$c\geq\sqrt2$, the bound $a(B)\leq c$ gives
\eqref{eq:inverse-support-global} as in that lemma; for smaller
$c$ the local bound already implies it.  Taking direct sums and
undoing the permutations completes the proof.
\end{proof}

\subsection{Retention of a prescribed maximum simplex}

The next reduction records why the matrix estimates control the
entire body, including points beyond the selected near-vertex
comparison simplex.

\begin{lemma}[From near-vertex points to inverse minors]
\label{lem:inverse-geometric-reduction}
Let $P$ have $n$ vertices, let $K\subset P$ be a convex body, and
let $S$ be a prescribed maximum-volume simplex in $K$.  Let $W$ be
the matrix of its vertex coordinates in the barycentric system of
$P$.  Suppose $K$ contains $n$ points whose nonnegative
column-stochastic barycentric matrix $Q$ satisfies
$Q_{ii}\geq1-\rho$.  For $0\leq\rho<1/2$, put
\[
 B=W^{-1},\qquad A=BQ,\qquad c=(1-2\rho)^{-1}.
\]
Then every square minor of $A$ has absolute value at most one,
and $M_k(B)\leq c^k$ for every $k$.  Moreover,
\begin{equation}
 E(K,S)\leq n\min\{1,a(B)\}.
 \label{eq:inverse-all-body-excess}
\end{equation}
The latter conclusion applies to all of $K$ and uses the original
centroid of the prescribed $S$.
\end{lemma}

\begin{proof}
Replacing any selected set of vertices of $S$ by selected columns
of $Q$ gives a simplex in $K$.  The determinant ratio is the
corresponding square minor of $A$, up to sign.  Maximality of $S$
therefore gives the all-minor bound.  Since
$\|Q-\Id_n\|_1\leq2\rho$, the Neumann series gives
$\|Q^{-1}\|_1\leq c$.  From $B=AQ^{-1}$, Cauchy--Binet bounds
each $k$-minor of $B$ by the sum of absolute $k$-minors of
$Q^{-1}$ with a fixed column set.  Expanding their determinants
over injections bounds that sum by the product of the corresponding
$k$ absolute column sums, at most $c^k$.

For any $x\in K$, its $P$-barycentric vector $z$ is nonnegative and
sums to one.  Its $S$-barycentric vector is $Bz$, whose every entry
is at least $\min B$.  Single-vertex replacement maximality also
gives $|\alpha_i^S(x)|\leq1$.  The original-centroid dilation
criterion \eqref{eq:dilationcoords} is
$\alpha_i^S(x)\geq-E/n$.  Combining these two lower bounds for
every $x\in K$ proves \eqref{eq:inverse-all-body-excess}.
\end{proof}

\begin{theorem}[Edge-supported prescribed simplex]
\label{thm:inverse-edge-simplex}
Under the hypotheses of Lemma~\ref{lem:inverse-geometric-reduction},
suppose every vertex of the originally prescribed maximum simplex
$S$ lies on an edge of $P$, allowing vertices of $P$ themselves.
For every $\rho\geq0$ for which the stated near-vertex hypothesis
holds,
\begin{equation}
 E(K,S)\leq(4+2\sqrt2)n\rho.
 \label{eq:inverse-edge-geometric}
\end{equation}
No bound on the number of vertices of $K$ is required.
\end{theorem}

\begin{proof}
The support assumption says every column of $W$ has at most two
nonzero entries.  Put
\[
 \rho_0=\frac{1-1/\sqrt2}{2}.
\]
For $0\leq\rho<\rho_0$, the reduction has
$c=(1-2\rho)^{-1}<\sqrt2$.  By
Theorem~\ref{thm:inverse-support-two} and
\eqref{eq:inverse-all-body-excess},
\[
 \frac{E(K,S)}n\leq c^2-1
 =\frac{4\rho(1-\rho)}{(1-2\rho)^2}
 \leq(4+2\sqrt2)\rho.
\]
Indeed $4(1-\rho)/(1-2\rho)^2$ increases on this interval and
equals $4+2\sqrt2$ at $\rho_0$.  For $\rho\geq\rho_0$, universal
single-replacement maximality gives $E(K,S)\leq n$, and
$1/\rho_0=4+2\sqrt2$.  This proves the entire large-$\rho$ branch,
including $\rho\geq1/2$ without using the Neumann reduction.
\end{proof}

\begin{theorem}[Every prescribed maximizer of an edge-supported body]
\label{thm:inverse-edge-body}
Let $K\subset P$ be a full-dimensional polytope whose every vertex
lies on an edge of $P$.  Suppose $K$ contains the near-vertex points
in Lemma~\ref{lem:inverse-geometric-reduction}.  Then
\eqref{eq:inverse-edge-geometric} holds for every prescribed
maximum-volume simplex $S$ of $K$, even if its own vertices are
not edge-supported.
\end{theorem}

\begin{proof}
Let $v_1,\ldots,v_N$ be the vertices of $K$.  Write every vertex
$s_j$ of the prescribed $S$ as a convex combination
$s_j=\sum_\ell\lambda_{\ell j}v_\ell$.  Independently sample
$T_j=v_\ell$ with probabilities $\lambda_{\ell j}$.  Use homogeneous
affine coordinates, in which the determinant of the $n$ columns
is proportional to oriented simplex volume, and order $S$ so that
$\det S=D>0$.  Determinant multilinearity gives
\[
 \mathbb E[\det T]=\det S=D.
\]
Maximality gives $|\det T|\leq D$ for every sampled tuple.  Hence
every positive-probability tuple has $\det T=D$: it is
nondegenerate, has the same orientation, and is a maximum vertex
simplex.

Fix $x\in K$ and an index $j$.  Cramer's rule and multilinearity
in all unreplaced columns give
\begin{align}
 \alpha_j^S(x)
 &=\frac{\det(S\text{ with column }j\text{ replaced by }x)}D
 \nonumber\\
 &=\frac{\mathbb E[\det(T\text{ with column }j
                              \text{ replaced by }x)]}D
  =\mathbb E[\alpha_j^T(x)].
 \label{eq:inverse-maximum-averaging}
\end{align}
Each sampled $T$ is edge-supported.  The coordinate bound in the
proof of Theorem~\ref{thm:inverse-edge-simplex} therefore gives
$\alpha_j^T(x)\geq-(4+2\sqrt2)\rho$ for all $x,j$.
Averaging proves the identical coordinate bound for the original
$S$.  The dilation criterion then yields
\eqref{eq:inverse-edge-geometric} with $S$ and its own centroid.
No sampled centroid is substituted.
\end{proof}

\begin{corollary}[Compact bodies supported on the one-skeleton]
\label{cor:inverse-extreme-skeleton}
Theorem~\ref{thm:inverse-edge-body} also applies when $K$ is a compact
convex body and every extreme point of $K$ lies on the finite
one-skeleton of $P$.
\end{corollary}

\begin{proof}
For every edge $F$ of $P$, $K\cap F$ is empty, a point, or a closed
interval.  A point in the relative interior of a nontrivial such
interval is not extreme in $K$.  Thus each edge contains at most
two extreme points of $K$, its interval endpoints.  There are only
finitely many edges, so $K$ has finitely many extreme points and,
by the finite-dimensional compact convex hull theorem, is their
convex hull.  It is therefore a polytope to which
Theorem~\ref{thm:inverse-edge-body} applies.  In particular, the same
conclusion holds when $K$ is the convex hull of finitely many closed
segments on the edges of $P$.
\end{proof}

\subsection{A conditional corollary for the actual defect}

\begin{corollary}[Actual-defect retention with compatible edge support]
\label{cor:inverse-actual-defect}
Retain the original actual invariant
\begin{equation}
 e(K)=\left(\frac d{d+1}\right)^d
       \frac{R(\pyr K)}{R(K)}-1-\frac1{d+1},
 \qquad
 R(L)=\frac{\vol(\Pi L)}{\vol(L)^{\dim L-1}},
 \label{eq:inverse-original-defect}
\end{equation}
and put $H_d=3d^2(d+1)^3(d+2)$, $n=d+1$, $d\geq3$.
On the branch $0\leq e(K)\leq H_d^{-1}$, assume additionally that
at least one enclosing simplex $P$ admitted by the actual-defect,
projection, and cap chain has every extreme point of $K$ on its
one-skeleton.  Then for every originally prescribed maximum simplex
$S$,
\begin{equation}
 E(K,S)\leq(4+2\sqrt2)(d+1)H_d^{1/(d-1)}
                         e(K)^{1/(d-1)}.
 \label{eq:inverse-actual-edge}
\end{equation}
In particular,
\begin{equation}
 E(K,S)\leq(768+384\sqrt2)d\,e(K)^{1/(d-1)}.
 \label{eq:inverse-actual-edge-coarse}
\end{equation}
For $e(K)>H_d^{-1}$, both bounds hold without the compatibility or
edge-support hypothesis.
\end{corollary}

\begin{proof}
On the stated small-defect branch, the retained actual chain
supplies the compatible enclosing $P$, relative projection deficit
at most $\delta=H_de(K)$, and near-vertex radius
$\rho=\delta^{1/(d-1)}$.  The extreme-point assumption makes $K$
a polytope by Corollary~\ref{cor:inverse-extreme-skeleton}.
Substitution into Theorem~\ref{thm:inverse-edge-body} gives
\eqref{eq:inverse-actual-edge}; the zero-defect endpoint is included
by the zero-radius cap conclusion.  The retained estimate
$H_d^{1/(d-1)}\leq144$, together with $d+1\leq4d/3$, gives
\eqref{eq:inverse-actual-edge-coarse}, since
$192(4+2\sqrt2)=768+384\sqrt2$.
For completeness, the numerical bound on $H_d$ follows from
$H_3=8640<144^2$ and
$H_{d+1}/H_d\leq25/6<144$ for $d\geq3$, by induction.
If $e(K)>H_d^{-1}$, then $H_d^{1/(d-1)}e(K)^{1/(d-1)}>1$,
so the universal $E(K,S)\leq d+1$ proves
\eqref{eq:inverse-actual-edge} directly; the coarser bound follows.
\end{proof}

\begin{remark}[The compatibility quantifier]
\label{rem:inverse-compatible-P}
The hypothesis refers to a $P$ produced by the actual-defect chain.
The existence of some unrelated enclosing simplex with edge-supported
extreme points is insufficient.  No result here shows that the
chain's $P$ can always be chosen with this property.  If only one
prescribed $S$ is to be retained, it is enough that this $S$ be
edge-supported relative to an admissible $P$, by
Theorem~\ref{thm:inverse-edge-simplex}.  The body-level hypothesis is
what gives the conclusion simultaneously for every prescribed
maximum simplex without inspecting each support.  The exact
coefficient in \eqref{eq:inverse-actual-edge} is asymptotic to
$(4+2\sqrt2)d$, since $H_d^{1/(d-1)}\to1$.  These are conditional
corollaries of the original invariant and the retained upstream
chain, not an unconditional theorem for all convex bodies.
\end{remark}

\subsection{Order three and an exact golden-ratio obstruction}

\begin{theorem}[An order-three two-minor inequality]
\label{thm:inverse-order-three}
If $n=3$, then
\begin{equation}
 1-\min_{i,j}B_{ij}\leq M_2(B).
 \label{eq:inverse-order-three}
\end{equation}
Consequently, if $M_1(B)\leq c$ and $M_2(B)\leq c^2$, then
\begin{equation}
 a(B)\leq\min\{c,c^2-1\}\leq\varphi^2(c-1),
 \qquad \varphi=\frac{1+\sqrt5}{2}.
 \label{eq:inverse-order-three-linear}
\end{equation}
The coefficient $\varphi^2$ is optimal.  The same conclusions hold
for direct sums of blocks of orders at most three, after independent
row and column permutations.
\end{theorem}

\begin{proof}
Permutations preserve stochasticity and the quantities in
\eqref{eq:inverse-order-three}.  First arrange
$D=\det W>0$.  Independent even row and column permutations are
transitive on the three indices, so any chosen entry of $B$ can be
placed at position $(3,1)$ while preserving $D>0$.  Write
\[
 W=\begin{pmatrix}
 u&v&z\\r&s&t\\1-u-r&1-v-s&1-z-t
 \end{pmatrix}.
\]
If $C_{13}$ is the cofactor at row $1$, column $3$, direct
expansion gives
\begin{align*}
 D-C_{13}
 &=(s-r)(1-z)+t(v-u)\\
 &=(1-z-t)(s-r)+t[(v+s)-(u+r)]+z\cdot0.
\end{align*}
This is a convex combination, with the third-column weights
$1-z-t,t,z$, of the values
$s-r$, $W_{31}-W_{32}$, and $0$.  Each is at most
$w_*:=\max_{i,j}W_{ij}$.  Since $B_{31}=C_{13}/D$,
\[
 1-B_{31}\leq\frac{w_*}{D}.
\]
Jacobi's complementary-minor identity gives $M_2(B)=w_*/D$.
Applying the argument to every entry proves
\eqref{eq:inverse-order-three}.  Thus $a(B)\leq c^2-1$, and the
entry bound gives $a(B)\leq c$.  For $c>1$, the upper envelope
of $\min\{c,c^2-1\}/(c-1)$ has maximum $\varphi^2$ at the crossing
$c=\varphi$; $c=1$ follows directly.

A one-by-one block is $[1]$.  In order two, every
column-stochastic inverse can be written
\[
 \begin{pmatrix}x&y\\1-x&1-y\end{pmatrix}.
\]
Each negative entry has a positive entry in its column larger in
absolute value by one.  Hence $a(B)\leq M_1(B)-1\leq c-1$ in
that order.  Applying the appropriate bound in each block proves
the direct-sum assertion.  Optimality follows from the next exact
example.
\end{proof}

\begin{proposition}[An all-order golden-ratio extremizer]
\label{prop:inverse-golden-example}
Let $p=\varphi^{-1}$ and set
\[
 W=\begin{pmatrix}
 0&p&1/2\\p&0&1/2\\1-p&1-p&0
 \end{pmatrix}.
\]
Then $W$ is nonnegative and column-stochastic,
$\det W=p(1-p)=\varphi^{-3}$, and
\[
 B=W^{-1}=\begin{pmatrix}
 -\varphi/2&\varphi/2&\varphi^2/2\\
 \varphi/2&-\varphi/2&\varphi^2/2\\
 1&1&-\varphi
 \end{pmatrix}.
\]
Its all-order root bound is exactly $c=\varphi$, and
\begin{equation}
 M_1(B)=\varphi,\quad M_2(B)=\varphi^2,\quad
 M_3(B)=\varphi^3,\quad
 a(B)=\varphi=c^2-1=\varphi^2(c-1).
 \label{eq:inverse-golden-equality}
\end{equation}
Thus any universal estimate $a(B)\leq C(c-1)$ valid for all
$c\geq1$ must have $C\geq\varphi^2$.
\end{proposition}

\begin{proof}
The stochasticity and displayed inverse follow by direct
multiplication using $\varphi^2=\varphi+1$.  The largest absolute
entry of $B$ is $\varphi$.  Jacobi gives
$M_2(B)=p/\det W=\varphi^2$, and
$M_3(B)=1/\det W=\varphi^3$.  The minimum entry is $-\varphi$,
which yields all equalities in \eqref{eq:inverse-golden-equality}.
\end{proof}

This is an exact obstruction to a smaller global constant.  It is
not a counterexample in the limit $c\downarrow1$.  The order-three
inequality \eqref{eq:inverse-order-three} is dimension-specific;
the corresponding proposed estimate with
$\max\{M_1(B),M_2(B)\}$ is false in order five.  In particular, an
unrestricted extension based only on first- and second-order minors
would require a new argument.

\subsection{Positive definiteness and a local permutation basin}

\begin{theorem}[Symmetric positive-definite case]
\label{thm:inverse-positive-definite}
If $W$ is symmetric and positive definite, then $M_1(B)\leq c$
alone implies
\begin{equation}
 a(B)\leq c-1.
 \label{eq:inverse-PSD}
\end{equation}
The coefficient one is sharp even under all-order minor bounds.
In the setting of Lemma~\ref{lem:inverse-geometric-reduction}, a
symmetric positive-definite prescribed $W$ gives the all-scale bound
\begin{equation}
 E(K,S)\leq4n\rho.
 \label{eq:inverse-PSD-geometric}
\end{equation}
\end{theorem}

\begin{proof}
A nonnegative column-stochastic matrix has spectral radius at most
one.  Symmetry and positive definiteness put all eigenvalues of $W$
in $(0,1]$.  Therefore $H=B-\Id_n$ is positive semidefinite and
\[
 0\leq H_{ii}=B_{ii}-1\leq c-1.
\]
Every two-coordinate positive-semidefinite restriction gives, for
$i\ne j$,
\[
 |B_{ij}|=|H_{ij}|\leq\sqrt{H_{ii}H_{jj}}\leq c-1.
\]
The diagonal entries are at least one, so this proves
\eqref{eq:inverse-PSD}.

For sharpness take
$W=\left(\begin{smallmatrix}1-s&s\\s&1-s\end{smallmatrix}\right)$
with $0<s<1/2$.  Its inverse has negative magnitude
$a=s/(1-2s)$, diagonal $1+a$, and determinant $1+2a$.
Since $(1+a)^2\geq1+2a$, its all-order root bound is $c=1+a$.
Thus equality holds in \eqref{eq:inverse-PSD}.

For the geometric claim, if $0\leq\rho<1/2$ then
\eqref{eq:inverse-all-body-excess} gives
$E(K,S)/n\leq\min\{1,c-1\}$ with $c=(1-2\rho)^{-1}$.
For $\rho\leq1/4$,
$c-1=2\rho/(1-2\rho)\leq4\rho$; for $\rho\geq1/4$ the
universal bound $E(K,S)/n\leq1\leq4\rho$ applies, also beyond
the Neumann range.
\end{proof}

The symmetry and positive-definiteness assumption is
coordinate-dependent.  It must be verified and does not follow
from positivity of the entries of $W$.

\begin{lemma}[A local permutation basin]
\label{lem:inverse-local-basin}
Suppose independent row and column permutations put $W$ in a form
with
\[
 s=\max_j(1-W_{jj})<\frac14.
\]
Then $M_1(B)\leq c$ implies
\begin{equation}
 a(B)\leq\frac{c-1}{1-4s}.
 \label{eq:inverse-local-basin}
\end{equation}
In particular, $s\leq1/8$ gives $a(B)\leq2(c-1)$.
\end{lemma}

\begin{proof}
Put $H=\Id_n-W$.  Its induced column norm is
$\|H\|_1=2s$, and its maximum entry modulus is at most $s$.
For $k\geq1$,
\[
 \max_{i,j}|(H^k)_{ij}|\leq s(2s)^{k-1}.
\]
The Neumann series therefore gives $a(B)\leq s/(1-2s)$.
At an index $j$ with $H_{jj}=s$, it also gives
\[
 B_{jj}\geq1+s-\frac{2s^2}{1-2s}
 =1+\frac{s(1-4s)}{1-2s}.
\]
Since $B_{jj}\leq c$, combining these inequalities proves
\eqref{eq:inverse-local-basin}.  The case $s=0$ is immediate.
\end{proof}

This lemma does not assert the existence of a suitable permutation.
A dimension-free argument placing arbitrary instances in such a
basin would itself constitute a substantive missing step.

\subsection{Additional all-order reductions}

The following results preserve useful information beyond the
support-two theorem.  Here the hypothesis is again
$M_k(B)\leq c^k$ for every relevant order.

\begin{proposition}[Localized minor completion and occupancy]
\label{prop:inverse-occupancy}
Let $R,S\subset[n]$ be disjoint row-label sets, with
$|R|=k$ and $|S|=\ell$, and let $J\subset[n]$ have size $k$.
For each $s\in S$, independently draw $Y_s$ from the distribution
$W_{:,s}$.  Then
\begin{equation}
 |\det B_{R,J}|
 \leq c^{k+\ell}\mathbb P\bigl(
     Y_s\text{ are pairwise distinct and all outside }J\bigr).
 \label{eq:inverse-occupancy}
\end{equation}
In particular, for any $B_{ij}$ and $S\subset[n]\setminus\{i\}$,
\begin{align}
 |B_{ij}|&\leq c^{|S|+1}
  \mathbb P\bigl(Y_s\text{ are distinct and all different from }j\bigr)
  \nonumber\\
 &\leq c^{|S|+1}\prod_{s\in S}(1-W_{j,s}).
 \label{eq:inverse-occupancy-entry}
\end{align}
With $k=0$ and the empty determinant equal to one,
\begin{equation}
 1\leq c^{|S|}\mathbb P(Y_s\text{ are pairwise distinct}).
 \label{eq:inverse-injection}
\end{equation}
\end{proposition}

\begin{proof}
Order $R\cup S$ with $R$ first.  Form the $(k+\ell)$-square
matrix whose first $k$ columns are $B_{R\cup S,J}$ and whose
last $\ell$ columns are the coordinate vectors $e_s$ restricted to
$R\cup S$.  Its determinant is $\det B_{R,J}$ up to a fixed
ordering sign.  Substitute
\[
 e_s=\sum_h W_{h,s}B_{:,h}
\]
and expand multilinearly in the last $\ell$ columns.  A term
vanishes if its selected indices repeat or if one belongs to $J$.
Every remaining determinant is a $(k+\ell)$-minor of $B$, bounded
in magnitude by $c^{k+\ell}$.  The sum of the corresponding
nonnegative product weights is exactly the probability in
\eqref{eq:inverse-occupancy}.  Taking $k=1$ gives the first entry
bound.  Dropping the collision-free condition and using the
independent draws gives the product bound.  Taking $k=0$ gives
\eqref{eq:inverse-injection}.
\end{proof}

The collision-free term retains both the fixed minor and its
forbidden column set and can be stronger than the displayed product
bound.

\begin{proposition}[Regularization of a uniformly almost nonnegative inverse]
\label{prop:inverse-regularization}
Suppose $A$ is invertible, every square minor of $A$ has absolute
value at most one, $\ones^{\mathsf T}A=\ones^{\mathsf T}$, and
$V=A^{-1}$ satisfies
$V_{ij}\geq-\eta/n$ for some $\eta\geq0$.  With
$J=\ones\ones^{\mathsf T}$, define
\[
 W=\frac{V+(\eta/n)J}{1+\eta},\qquad
 Q=(1+\eta)\Id_n-\frac\eta n J,\qquad B=W^{-1}=AQ.
\]
Then $W$ is nonnegative and column-stochastic,
\begin{equation}
 |B_{ij}-A_{ij}|\leq2\eta,
 \qquad M_k(B)\leq(1+2\eta)^k\quad(1\leq k\leq n).
 \label{eq:inverse-regularization}
\end{equation}
Thus a universal bound $a(B)\leq C(c-1)$ in the stochastic
inverse-minor problem would imply
$-\min A\leq(2C+2)\eta$ under these hypotheses.
\end{proposition}

\begin{proof}
The assumed lower bound gives $W\geq0$, and
$\ones^{\mathsf T}V=\ones^{\mathsf T}$ gives its column sums.
Using the same column-sum identity for $A$ verifies $B=AQ$.
Entrywise,
\[
 B_{ij}-A_{ij}
 =\eta\left(A_{ij}-\frac1n\sum_h A_{ih}\right).
\]
The order-one minor bound on $A$ proves the first part of
\eqref{eq:inverse-regularization}.

For a fixed $k$-element column set, the only nonzero $k$-minors
of $Q$ have the same row set or exchange one index.  The principal
minor is
\[
 (1+\eta)^{k-1}(1+\eta-k\eta/n),
\]
and each of the $k(n-k)$ one-index-exchange minors has magnitude
$(\eta/n)(1+\eta)^{k-1}$.  Cauchy--Binet therefore gives
\begin{align*}
 |\det B_{I,J}|&\leq
 (1+\eta)^{k-1}
 \left[1+\eta+\frac{\eta k(n-k-1)}n\right]\\
 &\leq(1+\eta)^{k-1}[1+(k+1)\eta]
 \leq(1+2\eta)^k.
\end{align*}
For the last inequality, Bernoulli's inequality gives
\[
 \frac{(1+2\eta)^k}{(1+\eta)^{k-1}}
 =(1+\eta)\left(1+\frac\eta{1+\eta}\right)^k
 \geq1+(k+1)\eta.
\]
Applying the hypothetical stochastic bound with $c=1+2\eta$
and then the entrywise difference estimate gives
$-\min A\leq2C\eta+2\eta$, as asserted.
\end{proof}

This is a one-way reduction.  It does not establish that an
arbitrary stochastic instance admits the stronger regularized
representation in the opposite direction.

\begin{lemma}[Conditional stability from a dominant matching]
\label{lem:inverse-dominant-matching}
If, after a row permutation, $W_{ii}\geq1-\varepsilon$ for every
$i$, with $0\leq\varepsilon<1/2$, then
\begin{equation}
 \|W^{-1}-\Id_n\|_1\leq
 \frac{2\varepsilon}{1-2\varepsilon}.
 \label{eq:inverse-dominant-matching}
\end{equation}
In particular the same upper bound holds for the negative entry
magnitudes in the permuted inverse.
\end{lemma}

\begin{proof}
Nonnegativity and column stochasticity give
$\|\Id_n-W\|_1\leq2\varepsilon$.  The Neumann series yields
\[
 \|W^{-1}-\Id_n\|_1
 \leq\frac{\|\Id_n-W\|_1}{1-\|\Id_n-W\|_1}
 \leq\frac{2\varepsilon}{1-2\varepsilon}.
\]
Every negative entry is bounded by this norm, proving the claim.
\end{proof}

A dimension-free dominant-matching theorem of the form
$W_{i,\sigma(i)}\geq1-O(c-1)$ would settle the target through
this lemma.  No such matching theorem is established here.

\subsection{Limitations of the structural results}

The support-two proof uses the unicyclic graph decomposition
essentially.  A column supported on three or more rows does not
reduce to the same path recursion, and no reduction of arbitrary
$W$ to this class is supplied.  The all-order conjecture
$a(B)\leq c^2-1$ remains open in general; its analogue based only
on first- and second-order bounds is false in order five.
Neither an unrestricted dimension-free inverse-minor theorem nor
a counterexample to it follows from this section.  The special
central-$Q_t$ formulation from Section~\ref{sec:projection} and
the unrestricted prescribed-maximum linear dimension bound also
remain open.  In particular, a matrix construction by itself is
not an actual-defect lower-bound body: the all-direction projection
and enclosing-simplex requirements must still be verified.
